\SourcePage{113}
\SGAchapter{I·附录 II}{宇宙（N. Bourbaki）}{II. Appendice : Univers (par N. Bourbaki)}
\footnotetext[3]{经 N. Bourbaki 同意，我们在这里刊出他的秘密文稿。本文中的引文均指他那部博学的著作。}
\SGAsection{1}{宇宙的定义及初步性质}
\begin{SGAresult}{定义}{1}
\Term{set}{集合} $U$ 称为一个\Term{universe}{宇宙}，如果它满足下列条件：
\begin{enumerate}
\item[(U.I)] 若 $x\in U$ 且 $y\in x$，则 $y\in U$；
\item[(U.II)] 若 $x,y\in U$，则 $\{x,y\}\in U$；
\item[(U.III)] 若 $x\in U$，则 $\mathcal P(x)\in U$；
\item[(U.IV)] 若 $(x_\alpha)_{\alpha\in I}$ 是 $U$ 的一族元素，且 $I\in U$，则并集 $\bigcup_{\alpha\in I}x_\alpha$ 属于 $U$；
\item[(U.?)] 若 $x,y\in U$，则\Term{ordered-pair}{有序对} $(x,y)\in U$。
\end{enumerate}
\end{SGAresult}
\textbf{附记。}我想，既然已经决定在以后的版本中按照 Kuratowski 的方式定义\Term{ordered-pair}{有序对}，即 $(x,y)=\{\{x,y\},\{x\}\}$，条件 (U.?) 就是多余的，因为它由 (U.II) 推出。

\textbf{例。}1）空集是一个\Term{universe}{宇宙}，记为 $U_\circ$。

2）考虑用四个\Term{signature}{符号}“$\{$”“$\}$”“$,$”和“$\varnothing$”组成的\Term{nonempty}{非空}\Term{finite}{有限}词（参见《\Term{algebra}{代数}》I）。按词长 $n$ 归纳定义有意义的词：

(a) 当 $n=1$ 时，只有词 $\varnothing$ 有意义；

(b) 长为 $n$ 的词 $A$ 有意义，当且仅当存在 $p$ 个彼此不同、长度小于 $n$ 的有意义的词 $A_1,\ldots,A_p$（$p\geq1$），使得
\[
 A=\{A_1,A_2,\ldots,A_p\}.
\]
例如，$\varnothing$、$\{\varnothing\}$、$\{\{\varnothing\},\{\{\varnothing\}\},\{\varnothing,\{\varnothing\}\}\}$ 都是有意义的词。按长度归纳，显然每一个有意义的词都表示\Term{set-theory}{集合论}中的一个项；例如，若 $A=\{A_1,A_2,\ldots,A_p\}$，则 $A$ 表示以 $A_1,A_2,\ldots,A_p$ 为元素的\Term{set}{集合}；这些项显然是\Term{finite}{有限}集。令 $U_1$ 为这样得到的\Term{set}{集合}全体所成的\Term{set}{集合}。容易验证，$U_1$ 满足定义 1 的条件 (U.I)、(U.II)、(U.III) 和 (U.IV)（却不满足那个愚蠢的 (U.?)；只要肯把\Term{ordered-pair}{有序对}弄妥，这也不要紧）。所以 $U_1$ 是一个\Term{universe}{宇宙}。注意，$U_1$ 的元素都是\Term{finite}{有限}集，而 $U_1$ 是可数的。

以下各命题中，$U$ 表示一个\Term{universe}{宇宙}。
\begin{SGAresult}{命题}{1}
若 $X\in U$ 且 $y\subset x$，则 $y\in U$。
\end{SGAresult}
\begin{proof}
由 (U.III)，$\mathcal P(x)\in U$；于是 $y\in\mathcal P(x)$，再由 (U.I) 得 $y\in U$。
\end{proof}
\begin{SGAresult}{推论}{}
若 $x\in U$，则 $x$ 的任一\Term{quotient-set}{商集} $y$ 都是 $U$ 的元素。
\end{SGAresult}
\begin{proof}
$y$ 是 $\mathcal P(x)$ 的\Term{subset}{子集}。由 (U.III) 及命题 1 得 $y\in U$。
\end{proof}

\SourcePage{114}
\begin{SGAresult}{命题}{2}
若 $x\in U$，则 $\{x\}\in U$。
\end{SGAresult}
\begin{proof}
在 (U.II) 中取 $y=x$ 即得。
\end{proof}
\begin{SGAresult}{命题}{3}
由 $U$ 的元素组成的任一\Term{ordered-pair}{有序对}、三元组、四元组，都是 $U$ 的元素。
\end{SGAresult}
\begin{proof}
由附记或 (U.?)，\Term{ordered-pair}{有序对}的情形成立。再由 $(x,y,z)=((x,y),z)$ 得三元组的情形，由 $(x,y,z,t)=((x,y,z),t)$ 得四元组的情形。
\end{proof}
\begin{SGAresult}{命题}{4}
若 $X,Y\in U$，则 $X\times Y\in U$。
\end{SGAresult}
\begin{proof}
对于 $x\in X$，$\{x\}\times Y$ 是族 $\bigl(\{(x,y)\}\bigr)_{y\in Y}$ 的并集；由 (U.I)、(U.II) 及命题 3，$\{(x,y)\}\in U$，故由 (U.IV) 得 $\{x\}\times Y\in U$。最后，$X\times Y$ 是族 $(\{x\}\times Y)_{x\in X}$ 的并集，再由 (U.IV)，它也是 $U$ 的元素。
\end{proof}
\begin{SGAresult}{推论}{1}
若 $X,Y,Z,\ldots$ 都是 $U$ 的元素，则在 $X,Y,Z,\ldots$ 上构造的阶梯中的所有\Term{set}{集合}，都是 $U$ 的元素（参见第四章，\S\,）。
\end{SGAresult}
\begin{proof}
相继应用命题 4 和 (U.III) 即得。
\end{proof}
\begin{SGAresult}{推论}{2}
若 $(x_\alpha)_{\alpha\in I}$ 是 $U$ 的一族元素且 $I\in U$，则\Term{sum-set}{和集} $\coprod_{\alpha\in I}x_\alpha$ 是 $U$ 的元素。
\end{SGAresult}
\begin{proof}
这个\Term{sum-set}{和集}是乘积 $\bigl(\bigcup_{\alpha\in I}x_\alpha\bigr)\times I$ 的\Term{subset}{子集}（第二章），而由 (U.IV) 及假设，这个乘积的两个因子都是 $U$ 的元素。应用命题 4 和 1 即得。
\end{proof}
\begin{SGAresult}{命题}{5}
若 $X$ 和 $Y$ 都是 $U$ 的元素，则 $X$ 与 $Y$ 之间的任一\Term{correspondence}{对应}（特别是 $X$ 到 $Y$ 的任一映射）都是 $U$ 的元素。
\end{SGAresult}
\begin{proof}
这样的\Term{correspondence}{对应} $C$ 是一个三元组 $(X,Y,\Gamma)$，其中 $\Gamma$ 是 $X\times Y$ 的\Term{subset}{子集}，即 $C$ 的图（第二章，\S\,）。由命题 4 和 1，$\Gamma\in U$，再由命题 3 得 $C\in U$。
\end{proof}
\begin{SGAresult}{命题}{6}
若 $X,Y\in U$，则任何由 $X$ 与 $Y$ 之间的\Term{correspondence}{对应}（特别是 $X$ 到 $Y$ 的映射）组成的\Term{set}{集合} $Z$ 都是 $U$ 的元素。
\end{SGAresult}
\begin{proof}
$Z$ 是 $\{X\}\times\{Y\}\times\mathcal P(X\times Y)$ 的\Term{subset}{子集}。由命题 2 及命题 4 的推论，这个乘积是 $U$ 的元素；故由命题 1，$Z\in U$。
\end{proof}
\begin{SGAresult}{推论}{}
若 $(x_\alpha)_{\alpha\in I}$ 是 $U$ 的一族元素，且 $I\in U$，则 $\prod_{\alpha\in I}x_\alpha\in U$。
\end{SGAresult}
\begin{proof}
这个乘积是由 $I$ 到 $\bigcup_{\alpha\in I}x_\alpha$ 的映射组成的一个\Term{set}{集合}，而后一个并集由 (U.IV) 属于 $U$。
\end{proof}

\SourcePage{115}
\begin{SGAresult}{命题}{7}
若 $X$ 是 $U$ 的\Term{subset}{子集}，且其\Term{cardinal}{基数}不大于 $U$ 中某个元素的\Term{cardinal}{基数}，则 $X$ 是 $U$ 的元素。
\end{SGAresult}
\begin{proof}
取 $I\in U$，使 $\operatorname{card}(X)\leq\operatorname{card}(I)$。存在由 $I$ 的某个\Term{subset}{子集} $I'$ 到 $X$ 的\Term{surjectivity}{满射} $i\mapsto x_i$。于是 $X$ 是族 $(\{x_i\})_{i\in I'}$ 的并集；由命题 1、命题 2 及 (U.IV)，这个并集属于 $U$。
\end{proof}
\begin{SGAresult}{推论}{}
若 $U$ \Term{nonempty}{非空}，则 $U$ 的任一\Term{finite}{有限}\Term{subset}{子集}都是 $U$ 的元素，并且 $U$ 有任意\Term{finite}{有限}\Term{cardinal}{基数}的元素。
\end{SGAresult}
反之，若 $U=\varnothing$，则 $\varnothing$ 是 $\varnothing$ 的\Term{finite}{有限}\Term{subset}{子集}，却不是 $\varnothing$ 的元素。
\begin{proof}
若 $U$ \Term{nonempty}{非空}，命题 2 表明某个单元素\Term{set}{集合} $x_\circ$ 属于 $U$。按 $n$ 归纳，令 $x_{n+1}=\mathcal P(x_n)$。由 (U.III)，$x_n\in U$，且 $\operatorname{card}(x_{n+1})=2^{\operatorname{card}(x_n)}$，所以 $U$ 有\Term{finite}{有限}\Term{cardinal}{基数}任意大的元素。
\end{proof}
\textbf{注记。}由命题 2 及命题 7 的推论，每个\Term{nonempty}{非空}\Term{universe}{宇宙} $U$ 都包含例 2 中的\Term{universe}{宇宙} $U_1$；事实上，由命题 1，$\varnothing\in U$。因此，$U_1$ 是所有\Term{nonempty}{非空}\Term{universe}{宇宙}的交集。更一般地：
\begin{SGAresult}{命题}{8}
若 $(U_\lambda)_{\lambda\in L}$ 是一个\Term{nonempty}{非空}的\Term{universe}{宇宙}族，则 $U=\bigcap_{\lambda\in L}U_\lambda$ 是一个\Term{universe}{宇宙}。
\end{SGAresult}
\begin{proof}
这立即由定义 1 得到。
\end{proof}

\SGAsection{2}{宇宙与结构种类}
设 $\mathscr E$ 是一种\Term{species-of-structure}{结构种类}；为使叙述具体而简便，假定每一个 $\mathscr E$ 类结构都定义在一个\Term{underlying-set}{基集}上。设 $(X,S)$ 是一个 $\mathscr E$ 类结构（$X$ 是\Term{underlying-set}{基集}，$S$ 是结构），$U$ 是一个\Term{universe}{宇宙}。若 $X\in U$，则 $X$ 上结构 $S$ 的所有构成\Term{object}{对象}都是 $U$ 的元素（由第 1 节命题 4 的推论 1 及 (U.I)），因而结构 $(X,S)$ 也是 $U$ 的元素。

现在假定 $\mathscr E$ 是带\Term{morphism}{态射}的\Term{species-of-structure}{结构种类}。若 $X,X'$ 是 $U$ 中配有 $\mathscr E$ 类结构的元素，则由 $X$ 到 $X'$ 的\Term{morphism}{态射}所成的\Term{set}{集合}仍是 $U$ 的元素（第 1 节命题 6）。

考虑第 1 节第 2 目例 c）中定义的\Term{category}{范畴} $(U\text{-}\mathscr E)$。由于 $(U\text{-}\mathscr E)$ 的\Term{object}{对象}和\Term{arrow}{箭头}都是 $U$ 的元素，$(U\text{-}\mathscr E)$ 是 $U$ 的两个\Term{subset}{子集}组成的\Term{ordered-pair}{有序对}，再配上由四个映射组成的四元组；一般说来，它不是 $U$ 的元素。记号 $X\in(U\text{-}\mathscr E)$ 表示 $X$ 是 $U$ 中配有 $\mathscr E$ 类结构的元素。

\Term{category}{范畴} $(U\text{-}\mathscr E)$ 对许多运算是\Term{stable}{稳定}的：

a) 若 $X\in(U\text{-}\mathscr E)$，$X'$ 是 $X$ 的一个容许\Term{induced-structure}{诱导结构}的\Term{subset}{子集}，则 $X'\in(U\text{-}\mathscr E)$。这由第 1 节命题 1 得到。

\SourcePage{116}
b) 若 $X\in(U\text{-}\mathscr E)$，$X''$ 是 $X$ 的一个容许\Term{quotient-structure}{商结构}的\Term{quotient-set}{商集}，则 $X''\in(U\text{-}\mathscr E)$（第 1 节命题 1 的推论）。

c) 若 $X,Y\in(U\text{-}\mathscr E)$，且 $X\times Y$ 容许\Term{product-structure}{乘积结构}，则 $X\times Y\in(U\text{-}\mathscr E)$（第 1 节命题 4）。更一般地，若 $(X_i)_{i\in I}$ 是 $(U\text{-}\mathscr E)$ 的一族元素，$I\in U$，且 $\prod_{i\in I}X_i$ 容许\Term{product-structure}{乘积结构}，则 $\prod_{i\in I}X_i\in(U\text{-}\mathscr E)$（命题 6 的推论）。对于\Term{sum-structure}{和结构}，有类似断言（命题 4 的推论 2）。

d) 设 $I$ 是\Term{preordered-set}{预序集}，$(X_i,f_{ij})_{i,j\in I}$ 是由配有 $\mathscr E$ 类结构的\Term{set}{集合}及其\Term{morphism}{态射}构成的一个\Term{projective-system}{投射系}（相应地，\Term{inductive-system}{归纳系}）；设 $L$ 是这个系在第三章意义下的\Term{limit}{极限}。若对每个 $i$ 都有 $X_i\in U$，且 $I\in U$，而 $L$ 容许\Term{projective-limit}{投射极限}结构（相应地，\Term{inductive-limit}{归纳极限}结构），则 $L\in(U\text{-}\mathscr E)$：事实上，$L$ 是 $\prod_{i\in I}X_i$ 的\Term{subset}{子集}（相应地，是 $\coprod_{i\in I}X_i$ 的\Term{quotient-set}{商集}），故由第 1 节，它是 $U$ 的元素。

再举几个更具体的例子：

* 1）设 $X,Y\in(U\text{-}\operatorname{Top})$ 是两个\Term{locally-compact-space}{局部紧空间}。则从 $X$ 到 $Y$ 的\Term{continuous-map}{连续映射}组成的\Term{set}{集合} $\mathscr C(X,Y)$，配以\Term{compact-convergence-topology}{紧致收敛拓扑}（《一般\Term{topology}{拓扑}》X），是 $(U\text{-}\operatorname{Top})$ 的一个元素。事实上，由第 1 节命题 6，$\mathscr C(X,Y)\in U$。*

* 2）设 $X,Y\in(U\text{-}\operatorname{Ab})$。则群 $\operatorname{Hom}_{\mathbf Z}(X,Y)$ 由第 1 节命题 6 是 $(U\text{-}\operatorname{Ab})$ 的元素。若进一步假定 $\mathbf Z\in U$，则\Term{tensor-product}{张量积} $X\otimes_{\mathbf Z}Y$ 是 $(U\text{-}\operatorname{Ab})$ 的元素；事实上，这个\Term{tensor-product}{张量积}是 $\mathbf Z^{(X\times Y)}$ 的商，而后者是 $\mathbf Z^{X\times Y}$ 的\Term{subset}{子集}，由第 1 节命题 4 和 6，$\mathbf Z^{X\times Y}$ 本身是 $U$ 的元素。*

* 3）设 $X\in(U\text{-}\operatorname{Unif})$ 是\Term{uniform-space}{一致空间}。则它的\Term{completion}{完备化} $\widehat X$ 是 $(U\text{-}\operatorname{Unif})$ 的元素。事实上，$\widehat X$ 是 $X$ 上 \Term{cauchy-filter}{Cauchy 滤子}的\Term{equivalence-class}{等价类}所成的\Term{set}{集合}；而 $X$ 上的一个\Term{filter}{滤子}是 $\mathcal P(\mathcal P(X))$ 的元素，所以\Term{filter}{滤子}的一个\Term{equivalence-class}{等价类}是 $\mathcal P(\mathcal P(\mathcal P(X)))$ 的元素，从而 $\widehat X$ 是 $\mathcal P(\mathcal P(\mathcal P(\mathcal P(X))))$ 的元素，故由 (U.III) 及 (U.I) 是 $U$ 的元素。*

\textbf{附记。}这番话的寓意是：Bourbaki 应当留心把自己的构造好好“\Term{canonical}{典范}化”。例如，在添入一个元素 $\infty$ 的构造中（Alexandroff \Term{compactification}{紧化}、\Term{projective}{射影}域），最好取 $\infty=\varnothing$（因为 $\varnothing$ 是每个\Term{nonempty}{非空}\Term{universe}{宇宙}的元素），再作 $\{\varnothing\}$ 与给定\Term{set}{集合}的\Term{sum-set}{和集}。当然，用来构造\Term{sum-set}{和集}的二元素\Term{set}{集合}也应当“\Term{canonical}{典范}化”；例如，$\{\varnothing,\{\varnothing\}\}$ 在我看来就是个好候选，因为它属于每个\Term{nonempty}{非空}\Term{universe}{宇宙}。

\SGAsection{3}{宇宙与范畴}
设 $U$ 是\Term{universe}{宇宙}，$C$ 是\Term{category}{范畴}。当 $\operatorname{Ob}(C)\in U$ 且 $\operatorname{Fl}(C)\in U$ 时，写作 $C\in U$。这个写法是合理的，因为 $C$ 是由 $\operatorname{Ob}(C)$、$\operatorname{Fl}(C)$ 及四个结构映射组成的六元组；所以，若 $\operatorname{Ob}(C)\in U$ 且 $\operatorname{Fl}(C)\in U$，
\SourcePage{117}
这四个映射也都是 $U$ 的元素（第 1 节命题 4 的推论 1），因而六元组 $C$ 也是如此（第 1 节命题 3，稍作通融）。若考虑\Term{category}{范畴}的\Term{species-of-structure}{结构种类}“$\operatorname{cat}$”，这也就是第 2 节的一个特例。用第 2 节的记号，关系 $C\in U$ 也写作 $C\in(U\text{-}\operatorname{cat})$。

注意，像 $(U\text{-}\operatorname{Ens})$、$(U\text{-}\operatorname{Ord})$、$(U\text{-}\operatorname{Top})$ 或 $(U\text{-}\operatorname{Gr})$ 这样的\Term{category}{范畴}，一般并不是 $U$ 的元素。

设 $C,D$ 是两个\Term{category}{范畴}，$U$ 是\Term{universe}{宇宙}，且 $C,D\in U$。若 $C'$ 是 $C$ 的\Term{subcategory}{子范畴}，则 $C'\in U$；\Term{product-category}{乘积范畴} $C\times D$、$C$ 与 $D$ 的\Term{sum-category}{和范畴}，以及\Term{opposite-category}{反范畴} $C^\circ,D^\circ$ 也都是 $U$ 的元素（参见第 2 节）。\Term{functor-category}{函子范畴} $E=\mathscr{H}\!om(C,D)$ 同样是 $U$ 的元素：事实上，$\operatorname{Ob}(E)$ 是由映射对 $\operatorname{Ob}(C)\to\operatorname{Ob}(D)$、$\operatorname{Fl}(C)\to\operatorname{Fl}(D)$ 组成的\Term{set}{集合}，故由第 1 节，$\operatorname{Ob}(E)\in U$；至于 $\operatorname{Fl}(E)$，它是\Term{functorial-morphism}{函子态射}所成的\Term{set}{集合}，即由映射 $\operatorname{Ob}(C)\to\operatorname{Fl}(D)$ 组成的\Term{set}{集合}，故由第 1 节命题 6，$\operatorname{Fl}(E)\in U$。

\begin{SGAresult}{命题}{9}
设 $\mathscr E$ 是带\Term{morphism}{态射}的\Term{species-of-structure}{结构种类}，$U,U'$ 是两个\Term{universe}{宇宙}，且 $U'\subset U$。则 $(U'\text{-}\mathscr E)$ 是 $(U\text{-}\mathscr E)$ 的\Term{full-subcategory}{满子范畴}。
\end{SGAresult}
\begin{proof}
若 $x,y$ 是 $(U'\text{-}\mathscr E)$ 的\Term{object}{对象}，则按定义它们是 $(U\text{-}\mathscr E)$ 的\Term{object}{对象}。而 $\operatorname{Hom}(x,y)$ 在这两个\Term{category}{范畴}中是同一个\Term{set}{集合}（参见第 1 节第 2 目例 c））。
\end{proof}
\textbf{附记。}注意，$U,U'$ 是\Term{universe}{宇宙}这个假设是多余的，所以命题 9 挪到第 1 节第 4 目去会更合适。不过部落要求把它放在这里。

\textbf{注记。}有时，\Term{species-of-structure}{结构种类} $\mathscr E$ 的公理蕴涵：配有 $\mathscr E$ 类结构的\Term{set}{集合}的\Term{cardinal}{基数}都以某个固定\Term{cardinal}{基数}为上界，记为 $\mathfrak{c}$（* 例如\Term{finite-group}{有限群}、\Term{finitely-generated}{有限生成}群、固定环 $A$ 上的\Term{finitely-generated}{有限生成}模，或 $A$ 上的\Term{finite-type}{有限型}\Term{algebra}{代数}的\Term{species-of-structure}{结构种类} *）。现在假定 $U'$ 中有一个元素 $z$，使 $\mathfrak{c}\leq\operatorname{card}(z)$。于是，对每个配有 $\mathscr E$ 类结构的\Term{set}{集合} $x$，$U'$ 中都有一个与 $x$ \Term{equipotent}{等势}的元素 $x'$，例如 $z$ 的某个\Term{subset}{子集}；通过\Term{transport-of-structure}{结构搬运}，由 $x$ 的结构给 $x'$ 配上结构，便得到 $(U'\text{-}\mathscr E)$ 的一个元素。因此，$(U'\text{-}\mathscr E)$ 到 $(U\text{-}\mathscr E)$ 的\Term{inclusion-functor}{包含函子}是\Term{essentially-surjective}{本质满射}（第 4 节第 2 目定义 2），从而是\Term{equivalence-of-categories}{范畴等价}（第 4 节第 3 目定理 1）。

\SGAsection{4}{宇宙公理}
第 2 节那些十分可喜的\Term{stable}{稳定}性结果，只有在能应用于第 1 节所述的两个小\Term{universe}{宇宙} $U_\circ,U_1$ 以外的情形时，才有意义。因此，我们在\Term{set-theory}{集合论}公理中添入下列公理：

(A.6) 对每个\Term{set}{集合} $x$，都存在一个\Term{universe}{宇宙} $U$，使 $x\in U$。

\SourcePage{118}
这个公理蕴涵：若 $(x_\alpha)_{\alpha\in I}$ 是一族\Term{set}{集合}，则存在一个\Term{universe}{宇宙} $U$，使对一切 $\alpha\in I$ 都有 $x_\alpha\in U$；事实上，只须对 $x=\bigcup_{\alpha\in I}x_\alpha$ 应用 (A.6)，再应用第 1 节命题 1 即可。

特别地，给定一个\Term{category}{范畴} $C$，存在一个\Term{universe}{宇宙} $U$，使 $C\in U$ 在第 3 节的意义下成立：只须对族 $(\operatorname{Ob}(C),\operatorname{Fl}(C))$ 应用上面的结论。这适用于形如 $(V\text{-}\mathscr E)$ 的\Term{category}{范畴}，其中 $\mathscr E$ 是带\Term{morphism}{态射}的\Term{species-of-structure}{结构种类}，$V$ 是一个\Term{set}{集合}，例如一个\Term{universe}{宇宙}；一般有 $V\ne U$。

例如，若 $V$ 是\Term{universe}{宇宙}，则 $\operatorname{Ob}(V\text{-}\operatorname{Ens})=V$ 且 $\operatorname{Fl}(V\text{-}\operatorname{Ens})\subset V$（第 1 节命题 5）。所以，使 $(V\text{-}\operatorname{Ens})\in U$ 的\Term{universe}{宇宙} $U$ 正是那些使 $V\in U$ 的\Term{universe}{宇宙}。而对于\Term{universe}{宇宙}，关系 $V\in V$ 不可能成立：否则，对 $V$ 的任一\Term{subset}{子集} $A$，都有 $A\in V$（第 1 节命题 1），于是 $\operatorname{card}(\mathcal P(V))\leq\operatorname{card}(V)$，这是不可能的。

\SGAsection{5}{宇宙与强不可达基数}
设 $U$ 是\Term{universe}{宇宙}。由定义 1 的条件 (U.I)，$U$ 的任一元素 $x$ 都是 $U$ 的\Term{subset}{子集}，故 $\operatorname{card}(x)\leq\operatorname{card}(U)$。由于 $U$ 的\Term{subset}{子集}的\Term{cardinal}{基数}构成一个\Term{well-ordered-set}{良序集}，\Term{cardinal}{基数}
\[
 \mathfrak{c}(U)=\sup_{x\in U}\operatorname{card}(x) \tag{1}
\]
存在。注意，对每个\Term{cardinal}{基数} $\mathfrak{c}<\mathfrak{c}(U)$，都有一个元素 $x\in U$，使 $\operatorname{card}(x)=\mathfrak{c}$；事实上，按定义存在 $y\in U$，使 $\mathfrak{c}\leq\operatorname{card}(y)\leq\mathfrak{c}(U)$，取 $y$ 的适当\Term{subset}{子集}作为 $x$ 即可。反过来，若 $x\in U$，则 $\operatorname{card}(x)<\mathfrak{c}(U)$；事实上，由 (U.III)，$\mathcal P(x)\in U$，所以 $2^{\operatorname{card}(x)}\leq\mathfrak{c}(U)$。因此，\Term{cardinal}{基数} $\mathfrak{c}(U)$ 具有下列性质：

1）若 $\mathfrak{c}$ 是小于 $\mathfrak{c}(U)$ 的\Term{cardinal}{基数}，则 $2^{\mathfrak{c}}<\mathfrak{c}(U)$；事实上，若 $x$ 是 $U$ 中\Term{cardinal}{基数}为 $\mathfrak{c}$ 的元素，则由 (U.III)，$\mathcal P(x)\in U$，所以 $2^{\mathfrak{c}}<\mathfrak{c}(U)$。

2）若 $(\mathfrak{c}_\lambda)_{\lambda\in I}$ 是一族\Term{cardinal}{基数}，对每个 $\lambda\in I$ 都有 $\mathfrak{c}_\lambda<\mathfrak{c}(U)$，且 $\operatorname{card}(I)<\mathfrak{c}(U)$，则\Term{cardinal}{基数}和 $\sum_{\lambda\in I}\mathfrak{c}_\lambda<\mathfrak{c}(U)$。事实上，取 $x_\lambda\in U$，使 $\operatorname{card}(x_\lambda)=\mathfrak{c}_\lambda$；把 $I$ 换成与它\Term{equipotent}{等势}的指标集后，可以假定 $I\in U$；于是各 $x_\lambda$ 的\Term{sum-set}{和集}是 $U$ 的元素（第 1 节命题 4 的推论 2），这就证明了断言。

作如下定义：
\begin{SGAresult}{定义}{2}
\Term{cardinal}{基数} $\mathfrak{d}$ 称为\Term{strongly-inaccessible-cardinal}{强不可达}的，如果：

(FI.1) 若 $\mathfrak{c}$ 是\Term{cardinal}{基数}且 $\mathfrak{c}<\mathfrak{d}$，则 $2^{\mathfrak{c}}<\mathfrak{d}$。

(FI.2) 若 $(\mathfrak{c}_\lambda)_{\lambda\in I}$ 是一族\Term{cardinal}{基数}，对每个 $\lambda\in I$ 都有 $\mathfrak{c}_\lambda<\mathfrak{d}$，且 $\operatorname{card}(I)<\mathfrak{d}$，则 $\sum_{\lambda\in I}\mathfrak{c}_\lambda<\mathfrak{d}$。
\end{SGAresult}
\textbf{例。}\Term{cardinal}{基数} $0$ 和可数无穷\Term{cardinal}{基数}都是\Term{strongly-inaccessible-cardinal}{强不可达}的。任何非零\Term{finite}{有限}\Term{cardinal}{基数}都不是\Term{strongly-inaccessible-cardinal}{强不可达}的。这样，我们刚刚证明了\Term{axiom-of-universes}{宇宙公理} (A.6) 蕴涵关系：
\SourcePage{119}

(A$'$.6) 每个\Term{cardinal}{基数}都严格小于某个\Term{strongly-inaccessible-cardinal}{强不可达基数}。

反过来：
\begin{SGAresult}{定理}{1}
关系 (A$'$.6) 蕴涵\Term{axiom-of-universes}{宇宙公理} (A.6)。
\end{SGAresult}
\begin{proof}
设 $A$ 是\Term{set}{集合}。须构造一个以 $A$ 为元素的\Term{universe}{宇宙} $U$。递归定义\Term{set}{集合}序列 $(A_n)_{n\geq0}$：
\[
 A_0=A,\qquad A_{n+1}=A_n\text{ 的元素的并集}=A_n\text{ 的元素所成的集合}。 \tag{2}
\]
令 $B=\bigcup_{n=0}^{\infty}A_n$。由 (A$'$.6)，取\Term{strongly-inaccessible-cardinal}{强不可达基数} $\mathfrak{c}$，使 $\operatorname{card}(B)<\mathfrak{c}$。

存在\Term{well-ordered-set}{良序集} $I$，使 $\operatorname{card}(I)=\mathfrak{c}$。将 $I$ 换成它的\Term{cardinal}{基数}为 $\mathfrak{c}$ 的最小\Term{initial-segment}{初始段}后，可以假定 $I$ 的每个异于 $I$ 的\Term{initial-segment}{初始段}，其\Term{cardinal}{基数}都小于 $\mathfrak{c}$。以 $\mathscr E$ 表示 $I$ 的最小元素。关于\Term{initial-segment}{初始段}的假设表明，$I$ 没有最大元素（否则可以把它去掉），所以 $I$ 的每个元素都有后继；对 $\alpha\in I$，以 $s(\alpha)$ 表示 $\alpha$ 的后继。

现在用\Term{transfinite-recursion}{超限递归}定义\Term{set}{集合}族 $(B_\alpha)_{\alpha\in I}$：
\[
 \begin{cases}
 B_{\mathscr E}=B,\\
 B_{s(\alpha)}=B_\alpha\cup\mathcal P(B_\alpha),\\
 B_\alpha=\displaystyle\bigcup_{\beta<\alpha}B_\beta,
       &\text{若 }\alpha\text{ 没有前驱。}
 \end{cases} \tag{3}
\]
令 $U=\bigcup_{\alpha\in I}B_\alpha$。我们将证明 $U$ 就是所求的\Term{universe}{宇宙}。

首先，$A\in U$，因为 $A=A_\circ$ 是 $B=B_{\mathscr E}$ 的\Term{subset}{子集}，故是 $\mathcal P(B_{\mathscr E})\subset B_{s(\mathscr E)}$ 的元素。

其次，注意对每个 $\alpha\in I$ 都有
\[
 \operatorname{card}(B_\alpha)<\mathfrak{c}. \tag{4}
\]
对 $\alpha$ 作\Term{transfinite-induction}{超限归纳}。$\alpha=\mathscr E$ 时由假设成立。由 (4) 可得
\[
 \operatorname{card}(B_{s(\alpha)})\leq
 \operatorname{card}(B_\alpha)+2^{\operatorname{card}(B_\alpha)}
 <\mathfrak{c}+\mathfrak{c}\quad\text{（由 (FI.1)）}=\mathfrak{c}
 \quad\text{（因为 }\mathfrak{c}\text{ 无限）。}
\]
最后，若 $\alpha$ 没有前驱，则按构造，$\beta<\alpha$ 所成的\Term{set}{集合} $I'$ 的\Term{cardinal}{基数}小于 $\mathfrak{c}$；故若对每个 $\beta<\alpha$ 都有 $\operatorname{card}(B_\beta)<\mathfrak{c}$，则
\[
 \operatorname{card}(B_\alpha)=\operatorname{card}\Bigl(\bigcup_{\beta\in I'}B_\beta\Bigr)
 \leq\sum_{\beta\in I'}\operatorname{card}(B_\beta)<\mathfrak{c}
 \quad\text{（由 (FI.2)）。}
\]
这就证明了 (4)。

现在证明
\[
 \operatorname{card}(x)<\mathfrak{c}\quad\text{对一切 }x\in U. \tag{5}
\]
只须证明：对每个 $\alpha\in I$，$x\in B_\alpha$ 蕴涵 $\operatorname{card}(x)<\mathfrak{c}$。再对 $\alpha$ 作\Term{transfinite-induction}{超限归纳}。当 $\alpha=\mathscr E$ 时成立，因为若 $x\in B_{\mathscr E}=B$，则某个 $n\geq0$ 满足 $x\in A_n$；由 (2)，$X\subset A_{n+1}$，故 $x\subset B$，且 $\operatorname{card}(x)\leq\operatorname{card}(B)<\mathfrak{c}$。$\alpha$ 没有前驱的情形显然。最后，若 $x\in B_\alpha$ 且 $\alpha=s(\beta)$，则或者 $x\in B_\beta$，这时断言“$\operatorname{card}(x)<\mathfrak{c}$”由归纳假设成立；
\SourcePage{120}
或者 $x\in\mathcal P(B_\beta)$，于是 $x\subset B_\beta$，这时断言“$\operatorname{card}(x)<\mathfrak{c}$”由 (4) 成立。

现在可以证明 $U$ 确实是\Term{universe}{宇宙}：

(U.I) 设 $x\in U$ 且 $y\in x$。须证明 $y\in U$。也就是说，须证明对每个 $\alpha\in I$ 都有
\[
 x\in B_\alpha\text{ 且 }y\in x\ \Longrightarrow\ y\in B_\alpha.
\]
仍对 $\alpha$ 作\Term{transfinite-induction}{超限归纳}。$\alpha=\mathscr E$ 时成立，因为 $x\in B_{\mathscr E}=B$ 蕴涵对某个 $n$ 有 $x\in A_n$；故若 $y\in x$，则 $y\in A_{n+1}$，从而 $y\in B=B_{\mathscr E}$。过渡到没有前驱的元素 $\alpha$ 是显然的。最后由 $\alpha$ 过渡到 $s(\alpha)$：若 $x\in B_{s(\alpha)}$ 且 $y\in x$，则或者 $x\in B_\alpha$，由归纳假设得 $y\in B_\alpha\subset B_{s(\alpha)}$；或者 $x\in\mathcal P(B_\alpha)$，仍有 $y\in B_\alpha\subset B_{s(\alpha)}$。

(U.II) 设 $x,y\in U$。由 (3)，族 $(B_\alpha)_{\alpha\in I}$ 递增，所以存在 $\alpha\in I$，使 $x,y\in B_\alpha$。于是 $\{x,y\}\in\mathcal P(B_\alpha)\subset B_{s(\alpha)}\subset U$。

(U.III) 设 $x\in U$。存在 $\alpha\in I$，使 $x\in B_\alpha$。所以只须证明对每个 $\alpha\in I$ 都有
\[
 x\in B_\alpha\ \Longrightarrow\ \mathcal P(x)\in B_{s(s(\alpha))}.
\]
仍对 $\alpha$ 作\Term{transfinite-induction}{超限归纳}。若 $\alpha=\mathscr E$，则对某个 $n$ 有 $x\in A_n$，故 $x$ 的任一\Term{subset}{子集} $y$ 都满足 $y\subset A_{n+1}\subset B=B_{\mathscr E}$；因此对每个 $y\in\mathcal P(x)$ 都有 $y\in\mathcal P(B_{\mathscr E})\subset B_{s(\mathscr E)}$，从而 $\mathcal P(x)\subset B_{s(\mathscr E)}$ 且 $\mathcal P(x)\in\mathcal P(B_{s(\mathscr E)})\subset B_{s(s(\mathscr E))}$。于是断言对 $\alpha=\mathscr E$ 成立。过渡到没有前驱的元素 $\alpha$ 是直接的，因为 $\beta<\alpha$ 此时蕴涵 $s(s(\beta))<\alpha$。最后由 $\alpha$ 过渡到 $s(\alpha)$；设 $x\in B_{s(\alpha)}$；若 $x\in B_\alpha$，则由归纳假设 $\mathcal P(x)\in B_{s(s(\alpha))}\subset B_{s(s(s(\alpha)))}$；若 $x\in\mathcal P(B_\alpha)$，则 $x\subset B_\alpha$，所以 $\mathcal P(x)\subset\mathcal P(B_\alpha)$ 且 $\mathcal P(x)\in\mathcal P(\mathcal P(B_\alpha))\in B_{s(s(s(\alpha)))}$。

(U.IV) 设 $(x_\lambda)_{\lambda\in K}$ 是 $U$ 的一族元素，且 $K\in U$。须证明并集 $x=\bigcup_{\lambda\in K}x_\lambda$ 是 $U$ 的元素。对每个 $\lambda\in K$ 选取 $\alpha(\lambda)\in I$，使 $x_\lambda\in B_{\alpha(\lambda)}$。证明各 $\alpha(\lambda)$ 所成的\Term{set}{集合} $\alpha(K)$ 在 $I$ 中有上界：否则便有
\[
 I=\bigcup_{\lambda\in K}[\mathscr E,\alpha(\lambda)],
\]
这与 (FI.2) 矛盾，因为 $\operatorname{card}(K)<\mathfrak{c}$，且对每个 $\lambda\in K$ 都有 $\operatorname{card}([\mathscr E,\alpha(\lambda)])<\mathfrak{c}$。因此取 $\alpha(K)$ 的一个上界 $\beta$；因族 $(B_\alpha)_{\alpha\in I}$ 递增，对每个 $\lambda\in K$ 都有 $x_\lambda\in B_\beta$。由 (U.I) 证明中所得的结论，故有 $x_\lambda\in B_\beta$，从而 $x=\bigcup_{\lambda\in K}x_\lambda\subset B_\beta$。由 (3)，$x\in B_{s(\beta)}$，所以 $x\in U$。
\end{proof}

\begin{SGAresult}{定义}{3}
设 $x,y$ 是两个\Term{set}{集合}，$n$ 是整数且 $n\geq0$。若存在序列 $(x_j)_{j=0,\ldots,n}$，使 $x_0=x$、$x_n=y$，且对 $j=0,\ldots,n-1$ 有 $x_{j+1}\in x_j$，就说 $y$ 是 $x$ 的一个 $n$ 阶\Term{set-component}{分量}。
\end{SGAresult}
因此 $x$ 是 $x$ 的唯一零阶\Term{set-component}{分量}。$x$ 的一阶（相应地，二阶）\Term{set-component}{分量}，就是 $x$ 的元素（相应地，$x$ 的元素的元素）。若存在 $n\geq0$，使 $y$ 是 $x$ 的一个 $n$ 阶\Term{set-component}{分量}，
\SourcePage{121}
就说 $y$ 是 $x$ 的一个\Term{set-component}{分量}。关系“$y$ 是 $x$ 的一个\Term{set-component}{分量}”是\Term{preorder}{预序关系}。由\Term{selection-union-schema}{选择—并集公理模式}（第二章，……），关系“$y$ 是 $x$ 的一个\Term{set-component}{分量}”关于 $y$ 是\Term{collectivizing-relation}{集体化}的，所以 $x$ 的所有\Term{set-component}{分量}构成一个\Term{set}{集合}。

\begin{SGAresult}{定义}{4}
设 $\mathfrak{c}$ 是\Term{cardinal}{基数}。若\Term{set}{集合} $x$ 的所有\Term{set-component}{分量}的\Term{cardinal}{基数}都 $\leq\mathfrak{c}$（相应地，都 $<\mathfrak{c}$，都\Term{finite}{有限}），则称 $x$ 为 $\mathfrak{c}$ 型（相应地，严格 $\mathfrak{c}$ 型，\Term{finite-type}{有限型}）。
\end{SGAresult}
\textbf{例。}\Term{universe}{宇宙} $U_1$（第 1 节例 2）的元素都是\Term{finite-type}{有限型}的。若 $\mathfrak{c}$ 是\Term{cardinal}{基数}，$x$ 为 $\mathfrak{c}$ 型（相应地，严格 $\mathfrak{c}$ 型，\Term{finite-type}{有限型}），则 $x$ 的每个\Term{set-component}{分量}和每个\Term{subset}{子集}也都是 $\mathfrak{c}$ 型（相应地，严格 $\mathfrak{c}$ 型，\Term{finite-type}{有限型}）；同样，$\mathcal P(x)$ 是 $2^{\mathfrak{c}}$ 型（相应地，严格 $2^{\mathfrak{c}}$ 型，\Term{finite-type}{有限型}）。若 $x$ 为 $\mathfrak{c}$ 型，且 $\mathfrak{c}\leq\mathfrak{c}'$，则 $x$ 为 $\mathfrak{c}'$ 型。

\begin{SGAresult}{引理}{}
设 $\mathfrak{c}$ 是不可数的\Term{strongly-inaccessible-cardinal}{强不可达基数}，$X$ 是严格 $\mathfrak{c}$ 型\Term{set}{集合}。则存在\Term{cardinal}{基数} $\mathfrak{d}<\mathfrak{c}$，使 $X$ 为 $\mathfrak{d}$ 型。
\end{SGAresult}
\begin{proof}
对每个整数 $n$，以 $\mathfrak{d}_n$ 表示 $X$ 的 $n$ 阶\Term{set-component}{分量}所成\Term{set}{集合} $X_n$ 的\Term{cardinal}{基数}。有 $\mathfrak{c}_0=1$、$\mathfrak{c}_1=\operatorname{card}(X)<\mathfrak{c}$；因 $X_n=\bigcup_{y\in X_{n-1}}y$，对 $n$ 归纳并使用 (FI.2)，便得对每个 $n$ 有 $\mathfrak{d}_n<\mathfrak{c}$。于是各 $\mathfrak{d}_n$ 都以\Term{cardinal}{基数} $\mathfrak{d}=\sum_{j\geq0}\mathfrak{d}_j$ 为上界，由 (FI.2) 及 $\mathfrak{c}$ 不可数的假设，$\mathfrak{d}<\mathfrak{c}$。若 $Y$ 是 $X$ 的一个 $n$ 阶\Term{set-component}{分量}，则 $\operatorname{card}(Y)\leq\operatorname{card}(X_{n+1})\leq\mathfrak{d}$。
\end{proof}
\begin{SGAresult}{命题}{10}
设 $U$ 是\Term{universe}{宇宙}，$\mathfrak{c}$ 是\Term{strongly-inaccessible-cardinal}{强不可达基数}。则 $U$ 中严格 $\mathfrak{c}$ 型的元素 $x$ 所成的\Term{set}{集合} $U'$ 是一个\Term{universe}{宇宙}。
\end{SGAresult}
\begin{proof}
若 $x,y\in U'$ 且 $z\in x$，则显然 $z\in U'$ 且 $\{x,y\}\in U'$。若 $x\in U'$，则 $\operatorname{card}(x)<\mathfrak{c}$，由 (FI.1)，$\operatorname{card}(\mathcal P(x))<\mathfrak{c}$；由于 $\mathcal P(x)$ 的一个\Term{set-component}{分量}，或者是 $\mathcal P(x)$，或者是 $x$ 的一个\Term{subset}{子集}，或者是 $x$ 的一个\Term{set-component}{分量}，故 $\mathcal P(x)\in U'$。最后，若 $(x_\alpha)_{\alpha\in I}$ 是 $U'$ 的一族元素，且 $I\in U'$，则由 (FI.2)，$\operatorname{card}(\bigcup_\alpha x_\alpha)<\mathfrak{c}$；因为 $\bigcup_\alpha x_\alpha$ 的任一正阶\Term{set-component}{分量}都是某个 $x_\alpha$ 的\Term{set-component}{分量}，故确有 $\bigcup_\alpha x_\alpha\in U'$。
\end{proof}
\textbf{关于\Term{cardinal}{基数} $\mathfrak{c}(U)$ 的注记。}设 $U$ 是\Term{universe}{宇宙}，$\mathfrak{c}(U)$ 是由 (1) 定义的\Term{cardinal}{基数}，即
\[
 \mathfrak{c}(U)=\sup_{x\in U}\operatorname{card}(x).
\]
显然 $\mathfrak{c}(U)\leq\operatorname{card}(U)$，因为如前所述，$x\in U$ 蕴涵 $x\subset U$。然而等式 $\mathfrak{c}(U)=\operatorname{card}(U)$ 并非总成立。例如，设 $(x_\alpha)_{\alpha\in I}$ 是一族不可数的\Term{signature}{符号}，并规定公理：
\[
 \text{“对一切 }\alpha\in I,\ x_\alpha=\{x_\alpha\}\text{”}
 \quad\text{和}\quad
 \text{“若 }\alpha\ne\beta,\text{ 则 }x_\alpha\ne x_\beta\text{”}
\]
（换言之，$x_\alpha$ 是以它自己为唯一元素的单元素\Term{set}{集合}）。像第 1 节例 2 那样，用\Term{signature}{符号} $\varnothing$、
\SourcePage{122}
$x_\alpha$、$\{$、$\}$ 等组成“有意义的词”。每个这样的词都表示一个\Term{finite}{有限}集。以 $U$ 表示这些\Term{set}{集合}所成的\Term{set}{集合}。容易验证定义 1 的各条件成立，故 $U$ 是\Term{universe}{宇宙}。由于有意义的词是某个\Term{cardinal}{基数}为 $\operatorname{card}(I)+4=\operatorname{card}(I)$ 的\Term{set}{集合}中元素组成的\Term{finite}{有限}序列，所以 $\operatorname{card}(U)=\operatorname{card}(I)$（第三章；事实上，只须 $\operatorname{card}(U)\geq\operatorname{card}(I)$ 就够了，而这是显然的）。另一方面，$\mathfrak{c}(U)$ 是可数\Term{cardinal}{基数}，故 $\mathfrak{c}(U)<\operatorname{card}(U)$。

严格不等式 $\mathfrak{c}(U)<\operatorname{card}(U)$ 的原因，在于这里引入了满足 $x_\alpha\in x_\alpha$ 的\Term{set}{集合} $x_\alpha$。若禁止这类怪物，就能对每个\Term{universe}{宇宙} $U$ 得到 $\operatorname{card}(U)=\mathfrak{c}(U)$，以及一些同样漂亮却“派不上任何用场”的结果。下一节我们就来做这件事。

\SGAsection{6}{Artin 集合与 Artin 宇宙}
\begin{SGAresult}{定义}{5}
若不存在无限序列 $(x_n)_{n\geq0}$，使 $x_\circ=x$ 且对每个 $n\geq0$ 都有 $x_{n+1}\in x_n$，则称\Term{set}{集合} $x$ 为 \Term{artinian-set}{Artin 集合}（ensemble artinien）。
\end{SGAresult}
\textbf{例。}\Term{set}{集合} $\varnothing$、$\{\varnothing\}$、$\{\varnothing,\{\varnothing\}\}$，以及更一般地第 1 节\Term{universe}{宇宙} $U_1$ 的各元素，都是 \Term{artinian-set}{Artin 集合}。

形象地说，若 $x$ 是 \Term{artinian-set}{Artin 集合}，从 $x$ 中取一个元素 $x_1$，再从 $x_1$ 中取一个元素 $x_2$，如此继续，这个过程必定终止，最终到达 $x$ 的一个为空集的\Term{set-component}{分量} $x_n$。换言之，\Term{artinian-set}{Artin 集合}是“从 $\varnothing$ 构造出来的”；以后将把这点说得更精确（参见定理 2 证明中的 (*)）。

\Term{artinian-set}{Artin 集合}显然具有下列性质：

(AR.I) 若 $x$ 是 \Term{artinian-set}{Artin 集合}，则 $x$ 的每个\Term{subset}{子集}及每个\Term{set-component}{分量}都是 \Term{artinian-set}{Artin 集合}。$x$ 为 \Term{artinian-set}{Artin 集合}，当且仅当 $x$ 的每个元素都是 \Term{artinian-set}{Artin 集合}。

(AR.II) 若 $x,y$ 是 \Term{artinian-set}{Artin 集合}，则 $\{x,y\}$ 是 \Term{artinian-set}{Artin 集合}。

(AR.III) 若 $x$ 是 \Term{artinian-set}{Artin 集合}，则 $\mathcal P(x)$ 是 \Term{artinian-set}{Artin 集合}。

(AR.IV) 任意 \Term{artinian-set}{Artin 集合}的并集都是 \Term{artinian-set}{Artin 集合}。

这些性质立即给出：
\begin{SGAresult}{命题}{11}
若 $U$ 是\Term{universe}{宇宙}，则 $U$ 的所有 \Term{artinian-set}{Artin 元素}所成的\Term{set}{集合}是一个\Term{universe}{宇宙}，而且必为 \Term{artinian-universe}{Artin 宇宙}。
\end{SGAresult}
\begin{SGAresult}{推论}{}
若 $x$ 是 \Term{artinian-set}{Artin 集合}，则存在 \Term{artinian-universe}{Artin 宇宙} $U$，使 $x\in U$。
\end{SGAresult}
\begin{proof}
由公理 (A.6)，$x$ 是某个\Term{universe}{宇宙} $V$ 的元素；取 $V$ 的所有 \Term{artinian-set}{Artin 元素}所成的\Term{set}{集合}为 $U$ 即可。
\end{proof}
下面的命题，比本节其余内容还更无用：
\begin{SGAresult}{命题}{12}
设 $A$ 是 \Term{artinian-set}{Artin 集合}。则：

a) 对每个 $x\in A$，有 $x\notin x$；
\SourcePage{123}

b) 若 $x,y\in A$，则 $x\in y$ 和 $y\in x$ 不可能同时成立；

c) 对每个 $x\in A$，在 $x$ 的\Term{set-component}{分量} $y,z$ 之间的关系“$y$ 是 $z$ 的一个\Term{set-component}{分量}”是\Term{order-relation}{序关系}；

d) 对 $A$ 的每个\Term{nonempty}{非空}元素 $x$，存在 $y\in x$，使 $x\cap y=\varnothing$。
\end{SGAresult}
\begin{proof}
a)（相应地，b)）的否定会推出与定义 5 相矛盾的无限序列 $(x_n)_{n\geq0}$，即 $(x,x,x,\ldots)$（相应地，$(x,y,x,y,x,y,\ldots)$）。c) 的否定是说：存在 $x\in A$、$x$ 的两个不同\Term{set-component}{分量} $y,z$，以及属于关系组成的序列
\[
 y\in y_1\in\cdots\in y_q\in z
 \quad\text{和}\quad
 z\in z_1\in\cdots\in z_r\in y;
\]
于是像 b) 中一样，可得与定义 5 矛盾的无限序列 $(x_n)_{n\geq0}$。最后，若 d) 不成立，则存在 $A$ 的一个\Term{nonempty}{非空}元素 $x$，使对每个 $y\in x$ 都有 $x\cap y\ne\varnothing$；令 $x_\circ=x$，取 $x$ 的一个元素为 $x_1$；因 $x\cap x_1\ne\varnothing$，取 $x\cap x_1$ 的一个元素为 $x_2$，等等。更形式地，用 $x_1\in x$、$x_{n+1}\in x_n\cap x$（$n\geq1$）递归定义 $x$ 中元素的无限序列 $(x_n)_{n\geq1}$；于是序列 $(x_n)_{n\geq0}$ 与定义 5 矛盾。
\end{proof}
\textbf{注记。}设 $B$ 是\Term{set}{集合}。$B$ 为 \Term{artinian-set}{Artin 集合}，当且仅当：任何由 $B$ 的\Term{set-component}{分量}的\Term{set}{集合}所组成的\Term{set}{集合} $A$，都满足命题 12 的条件 d)。必要性由 (AR.I) 及命题 12 得到。反过来，若 $B$ 不是 \Term{artinian-set}{Artin 集合}，则存在无限序列 $(x_n)_{n\geq0}$，使 $x_\circ=B$ 且对每个 $n\geq0$ 都有 $x_{n+1}\in x_n$；此时以 $X$ 表示各 $x_n$ 所成的\Term{set}{集合}，取只含\Term{set}{集合} $X$ 的单元素集作为 $A$；于是 $A$ 包含一个\Term{nonempty}{非空}元素 $X$，使对 $X$ 的每个元素 $y$ 都有 $y\cap X\ne\varnothing$（事实上，$y$ 具有形式 $x_n$，且 $x_{n+1}\in y\cap X$）。

\begin{SGAresult}{定理}{2}
设 $\mathfrak{c}$ 是无限\Term{cardinal}{基数}。则：

a) 关系“$x$ 是 $\mathfrak{c}$ 型（相应地，严格 $\mathfrak{c}$ 型）的 \Term{artinian-set}{Artin 集合}”关于 $x$ 是\Term{collectivizing-relation}{集体化}的；$\mathfrak{c}$ 型 \Term{artinian-set}{Artin 集合}所成的\Term{set}{集合} $A_{\mathfrak{c}}$ 的\Term{cardinal}{基数}为 $2^{\mathfrak{c}}$。

b) 若 $\mathfrak{c}$ \Term{strongly-inaccessible-cardinal}{强不可达}，则严格 $\mathfrak{c}$ 型 \Term{artinian-set}{Artin 集合}所成的\Term{set}{集合} $U_{\mathfrak{c}}$ 是一个\Term{cardinal}{基数}为 $\mathfrak{c}$ 的\Term{universe}{宇宙}；\Term{cardinal}{基数} $\mathfrak{c}(U_{\mathfrak{c}})=\sup_{x\in U_{\mathfrak{c}}}\operatorname{card}(x)$ 等于 $\mathfrak{c}$。

c) 若一个\Term{universe}{宇宙} $U$ 有一个\Term{cardinal}{基数}为 $\mathfrak{c}$ 的元素，则每个 $\mathfrak{c}$ 型 \Term{artinian-set}{Artin 集合}都属于 $U$（换言之，$A_{\mathfrak{c}}\subset U$）。

c$'$) 若\Term{universe}{宇宙} $U$ \Term{nonempty}{非空}，则每个\Term{finite-type}{有限型} \Term{artinian-set}{Artin 集合}都是 $U$ 的元素。
\end{SGAresult}
在证明定理 2 以前，先推出几个很能启发人的推论：
\begin{SGAresult}{推论}{1}
若\Term{universe}{宇宙} $U$ 是 \Term{artinian-set}{Artin 的}，则 $\operatorname{card}(U)$ \Term{strongly-inaccessible-cardinal}{强不可达}，并且 $U$ 就是严格 $\operatorname{card}(U)$ 型 \Term{artinian-set}{Artin 集合}所成的\Term{set}{集合}。
\end{SGAresult}
\begin{proof}
令 $\mathfrak{c}(U)=\sup_{x\in U}\operatorname{card}(x)$。这是一个\Term{strongly-inaccessible-cardinal}{强不可达基数}（第 5 节开头）。先假定它不可数；于是，对每个无限\Term{cardinal}{基数}
\SourcePage{124}
$\mathfrak{c}<\mathfrak{c}(U)$，由 c)，每个 $\mathfrak{c}$ 型 \Term{artinian-set}{Artin 集合}都是 $U$ 的元素；故由第 5 节的引理，每个严格 $\mathfrak{c}(U)$ 型 \Term{artinian-set}{Artin 集合}都是 $U$ 的元素。由 c$'$)，当 $\mathfrak{c}(U)$ 可数时，最后这个断言仍成立。反过来，由 (U.I)，$U$ 中的每个\Term{set}{集合}都是严格 $\mathfrak{c}(U)$ 型。所以 $U$ 就是 b) 中的\Term{universe}{宇宙} $U_{\mathfrak{c}(U)}$，再由 b)，$\mathfrak{c}(U)=\operatorname{card}(U)$。
\end{proof}
由推论 1，\Term{artinian-universe}{Artin 宇宙}由它的\Term{cardinal}{基数}（这个\Term{cardinal}{基数}还是\Term{strongly-inaccessible-cardinal}{强不可达}的）唯一确定。所以 \Term{artinian-universe}{Artin 宇宙}与\Term{strongly-inaccessible-cardinal}{强不可达基数}之间有一个“一一\Term{correspondence}{对应}”。特别地：
\begin{SGAresult}{推论}{2}
\Term{artinian-universe}{Artin 宇宙}之间的包含关系 $U\subset U'$ 是\Term{well-order}{良序关系}。
\end{SGAresult}
\begin{proof}
\Term{cardinal}{基数}之间的关系 $\mathfrak{c}\leq\mathfrak{c}'$ 是\Term{well-order}{良序关系}（第三章）；并且，用定理 2 b) 的记号，关系 $\mathfrak{c}\leq\mathfrak{c}'$ 与 $U_{\mathfrak{c}}\subset U_{\mathfrak{c}'}$ \Term{equivalence}{等价}。
\end{proof}
注意，定理 2 b) 给出了定理 1（第 5 节）的第二个证明。

现在开始证明定理 2。给定\Term{set}{集合} $A$，称任一\Term{finite}{有限}序列 $(x_j)_{j=0,\ldots,n}$ 为 $A$ 的一条链，如果 $x_0=A$ 且对 $j=0,\ldots,n-1$ 有 $x_{j+1}\in x_j$。由\Term{selection-union-schema}{选择—并集公理模式}，$A$ 的所有链构成一个\Term{set}{集合}，记为 $G(A)$。给定一条链 $X=(x_j)_{j=0,\ldots,n}$，形如 $(x_i)_{i=0,\ldots,q}$ 且 $q\leq n$ 的链称为比 $X$ 小的链；这便给 $G(A)$ 赋予\Term{ordered-set}{有序集}结构。$A$ 为 \Term{artinian-set}{Artin 集合}，当且仅当 $G(A)$ 是“Noether”\Term{ordered-set}{有序集}（即满足第三章第 6 节第 5 目中的那些\Term{equivalence}{等价}条件）。

我们将证明：

(*) 若 $A$ 是 \Term{artinian-set}{Artin 集合}，则它由\Term{ordered-set}{有序集} $G(A)$ 的\Term{isomorphism-class}{同构类}唯一确定。

事实上，给定\Term{ordered-set}{有序集} $G$ 及元素 $g\in G$，以 $S(g)$ 表示满足 $g<g'$，且 $g\leq h\leq g'$ 蕴涵 $h=g$ 或 $h=g'$ 的所有 $g'\in G$ 所成的\Term{set}{集合}（换言之，就是 $g$ 的所有“直接后继”所成的\Term{set}{集合}）。考虑把 $G(A)$ 的每条链 $X=(x_0,\ldots,x_n)$ 映到\Term{set}{集合} $x_n$ 的映射 $\theta$；于是对每个 $X\in G(A)$ 都有
\[
 \theta(X)=\{\theta(X')\mid X'\in S(X)\}.
\]
由于 $A$ 是 $G(A)$ 的最小元素 $X_0=(A)$ 在 $\theta$ 下的像，只须证明 $\theta$ 由\Term{ordered-set}{有序集} $G(A)$ 的\Term{isomorphism-class}{同构类}唯一确定即可。这由下面的引理得到：
\begin{SGAresult}{引理}{1}
设 $G$ 是 \Term{noetherian-ordered-set}{Noether 有序集}，$\varphi$ 是定义在 $G$ 上的映射，对每个 $g\in G$ 都有
\[
 \varphi(g)=\{\varphi(g')\mid g'\in S(g)\}. \tag{1}
\]
则 $\varphi$ 唯一确定。此外，若 $U$ 是包含一个与 $G$ \Term{equipotent}{等势}的元素的\Term{universe}{宇宙}，则 $\varphi$ 的值都落在 $U$ 中。
\end{SGAresult}
\SourcePage{125}
\begin{proof}
假设蕴涵：若 $g$ 是 $G$ 的极大元素，则 $\varphi(g)=\varnothing$。

设 $\varphi,\varphi'$ 是使 (1) 成立的两个映射；若 $\varphi\ne\varphi'$，则满足 $\varphi(g)\ne\varphi'(g)$ 的 $g\in G$ 所成的\Term{set}{集合}\Term{nonempty}{非空}，故因 $G$ 是 \Term{noetherian-ordered-set}{Noether 的}而有一个极大元素 $h$；这时对每个 $g>h$ 都有 $\varphi(g)=\varphi'(g)$，特别对每个“后继”$g\in S(h)$ 都如此；由 (1) 得 $\varphi(h)=\varphi'(h)$，矛盾。因此确有 $\varphi=\varphi'$。现在设\Term{universe}{宇宙} $U$ 包含一个与 $G$ \Term{equipotent}{等势}的元素 $x$；证明 $\varphi$ 的值落在 $U$ 中；否则，在满足 $\varphi(g)\notin U$ 的 $g\in G$ 中取一个极大元素 $h$；对每个 $g'\in S(g)$ 都有 $\varphi(g')\in U$，所以
\[
 \varphi(h)=\{\varphi(g')\mid g'\in S(h)\}
\]
是 $U$ 的\Term{subset}{子集}；而它的\Term{cardinal}{基数}不大于 $\operatorname{card}(G)$，因而不大于 $U$ 中某个元素的\Term{cardinal}{基数}，故 $\varphi(h)\in U$（第 1 节命题 7）。这个矛盾表明 $\varphi$ 的值落在 $U$ 中。
\end{proof}
接着证明定理 2 a)。只须证明不带括号的断言，因为括号中的相应断言立即由它推出。设 $\mathfrak{c}$ 是无限\Term{cardinal}{基数}。若 $A$ 是 $\mathfrak{c}$ 型\Term{set}{集合}，则
\[
 \operatorname{card}(G(A))\leq\mathfrak{c}. \tag{2}
\]
事实上，以 $A_n$ 表示 $A$ 的 $n$ 阶\Term{set-component}{分量}所成的\Term{set}{集合}，则 $\operatorname{card}(A_1)\leq\mathfrak{c}$，且 $\operatorname{card}(A_{n+1})\leq\mathfrak{c}\cdot\operatorname{card}(A_n)$；归纳得 $\operatorname{card}(A_n)\leq\mathfrak{c}^n$；因为 $\mathfrak{c}$ 无限，$\mathfrak{c}^n=\mathfrak{c}$（第三章）；于是 $\operatorname{card}(\bigcup_{n=0}^{\infty}A_n)\leq\operatorname{card}(\mathbf N)\cdot\mathfrak{c}=\mathfrak{c}$（第三章）；而 $G(A)$ 是由 $\bigcup_n A_n$ 中元素组成的\Term{finite}{有限}序列所成的\Term{set}{集合}，所以确有不等式 (2)（第三章）。

现在，若 $E$ 是\Term{cardinal}{基数}为 $\mathfrak{c}$ 的\Term{set}{集合}，则在 $E$ 的一个\Term{subset}{子集} $E'$ 上给定序结构，\Term{equivalence}{等价}于给定 $E\times E$ 中由满足 $x\in E'$、$y\in E'$ 且 $x\leq y$ 的 $(x,y)$ 组成的\Term{subset}{子集}。故由 (2)，各\Term{ordered-set}{有序集} $G(A)$（其中 $A$ 为 $\mathfrak{c}$ 型）的\Term{isomorphism-class}{同构类}构成一个\Term{set}{集合} $\mathscr F_{\mathfrak{c}}$，且
\[
 \operatorname{card}(\mathscr F_{\mathfrak{c}})\leq\operatorname{card}(\mathcal P(E\times E))=2^{\mathfrak{c}}.
\]
设 $\mathscr F'_{\mathfrak{c}}$ 是 $\mathscr F_{\mathfrak{c}}$ 中由具有最小元素的 \Term{noetherian-ordered-set}{Noether 有序集}的类所组成的\Term{subset}{子集}；若对每个 $G\in\mathscr F'_{\mathfrak{c}}$，令其\Term{correspondence}{对应}于引理 2 中函数 $\varphi$ 在 $G$ 的最小元素处的值，便得 $\mathscr F'_{\mathfrak{c}}$ 上的映射 $\theta$，它的像包含所有 $\mathfrak{c}$ 型 \Term{artinian-set}{Artin 集合}。因此，这些\Term{set}{集合}确实构成一个\Term{set}{集合} $A_{\mathfrak{c}}$，且
\[
 \operatorname{card}(A_{\mathfrak{c}})\leq\operatorname{card}(\mathscr F'_{\mathfrak{c}})
 \leq\operatorname{card}(\mathscr F_{\mathfrak{c}})\leq2^{\mathfrak{c}}.
\]
还须证明 $\operatorname{card}(A_{\mathfrak{c}})=2^{\mathfrak{c}}$。为此，只须证明存在一个\Term{cardinal}{基数}为 $\mathfrak{c}$ 的 $\mathfrak{c}$ 型 \Term{artinian-set}{Artin 集合} $B$，因为此时 $B$ 的各\Term{subset}{子集}都是 $A_{\mathfrak{c}}$ 的元素。它的存在性由下列引理得到：
\begin{SGAresult}{引理}{2}
对每个\Term{cardinal}{基数} $\mathfrak{c}$，都存在一个\Term{cardinal}{基数}为 $\mathfrak{c}$ 的 $\mathfrak{c}$ 型 \Term{artinian-set}{Artin 集合} $B$。
\end{SGAresult}
\begin{proof}
对 $\mathfrak{c}$ 作\Term{transfinite-induction}{超限归纳}。$\mathfrak{c}=0$ 时没有选择，只能取 $B=\varnothing$。若 $\mathfrak{c}$ 有前驱 $\mathfrak{c}'$，设 $B'$ 是\Term{cardinal}{基数}为 $\mathfrak{c}'$ 的 $\mathfrak{c}'$ 型 \Term{artinian-set}{Artin 集合}；有 $\mathfrak{c}\leq2^{\mathfrak{c}'}$，故存在 $\mathcal P(B')$ 的\Term{cardinal}{基数}为 $\mathfrak{c}$ 的\Term{subset}{子集} $B$；$B$ 的元素是 $B'$ 的\Term{subset}{子集}，故它们的\Term{cardinal}{基数}都 $\leq\mathfrak{c}'\leq\mathfrak{c}$；$B$ 的
\SourcePage{126}
更高阶\Term{set-component}{分量}都是 $B'$ 的\Term{set-component}{分量}，故它们的\Term{cardinal}{基数}也都 $\leq\mathfrak{c}'\leq\mathfrak{c}$。最后，若 $\mathfrak{c}$ 没有前驱，则对每个\Term{cardinal}{基数} $\mathfrak{c}_\lambda<\mathfrak{c}$，选取\Term{cardinal}{基数}为 $\mathfrak{c}_\lambda$ 的 $\mathfrak{c}_\lambda$ 型 \Term{artinian-set}{Artin 集合} $B_\lambda$；于是 $B=\bigcup_\lambda B_\lambda$ 就满足要求。这证明了引理 2，也完成了 a) 部分的证明。
\end{proof}
再看 b)。设 $\mathfrak{c}$ 是\Term{strongly-inaccessible-cardinal}{强不可达基数}。由 a) 已知，严格 $\mathfrak{c}$ 型 \Term{artinian-set}{Artin 集合}构成一个\Term{set}{集合} $U_{\mathfrak{c}}$。$U_{\mathfrak{c}}$ 是\Term{universe}{宇宙}，立即由 \Term{artinian-set}{Artin 集合}的性质 (AR.I) 至 (AR.IV)（本节开头）以及与第 5 节命题 10 中类似的\Term{cardinal}{基数}估计得到。把引理 2 应用于各个\Term{cardinal}{基数} $\mathfrak{d}<\mathfrak{c}$，得到 $\sup_{x\in U_{\mathfrak{c}}}\operatorname{card}(x)=\mathfrak{c}$。最后，为证明 $\operatorname{card}(U_{\mathfrak{c}})=\mathfrak{c}$，先设 $\mathfrak{c}$ 不可数；由第 5 节的引理，$U_{\mathfrak{c}}$ 是并集 $\bigcup_{\mathfrak{d}<\mathfrak{c}}A_{\mathfrak{d}}$，其中 $A_{\mathfrak{d}}$ 表示 $\mathfrak{d}$ 型 \Term{artinian-set}{Artin 集合}所成的\Term{set}{集合}；由 a)，$\operatorname{card}(A_{\mathfrak{d}})=2^{\mathfrak{d}}<\mathfrak{c}$；另一方面，所有小于 $\mathfrak{c}$ 的\Term{cardinal}{基数}所成的\Term{set}{集合}，其\Term{cardinal}{基数} $\leq\mathfrak{c}$（第三章）；故 $\operatorname{card}(U_{\mathfrak{c}})\leq\mathfrak{c}\cdot\mathfrak{c}=\mathfrak{c}$；又 $\operatorname{card}(U_{\mathfrak{c}})\geq\mathfrak{c}$，因为对每个 $\mathfrak{d}<\mathfrak{c}$ 都有 $\operatorname{card}(U_{\mathfrak{c}})\geq\operatorname{card}(A_{\mathfrak{d}})=2^{\mathfrak{d}}$。

$\mathfrak{c}=0$ 的情形\Term{trivial-bundle}{平凡}，只剩 $\mathfrak{c}$ 为可数无穷\Term{cardinal}{基数}的情形。这时 $U_{\mathfrak{c}}$ 是\Term{finite-type}{有限型} \Term{artinian-set}{Artin 集合}（即自身及其所有\Term{set-component}{分量}都\Term{finite}{有限}的\Term{set}{集合}）所成的\Term{set}{集合}，我们要使用一个漂亮的组合结果。
\begin{SGAresult}{引理}{3}
（D. König？）考虑由\Term{finite}{有限}集 $E_n$ 及映射 $f_n:E_{n+1}\to E_n$ 组成的两个无限序列 $(E_n)_{n\geq1}$、$(f_n)_{n\geq1}$。若不存在无限序列 $(x_n)_{n\geq1}$，使 $x_n\in E_n$ 且对每个 $n$ 有 $f_n(x_{n+1})=x_n$，则当 $n$ 充分大时 $E_n$ 为空。
\end{SGAresult}
换言之，若所有满足 $f_n(x_{n+1})=x_n$ 的序列 $(x_n)$ 都\Term{finite}{有限}，则它们的长度有界。这可以用\Term{projective-limit}{投射极限}来表述：\Term{nonempty}{非空}\Term{finite}{有限}集的\Term{projective-limit}{投射极限}\Term{nonempty}{非空}（参见《一般\Term{topology}{拓扑}》第一章，第 2 版，第 9 节第 6 目命题 8，2°）。
\begin{proof}
用反证法。若对任意大的 $n$ 都存在\Term{nonempty}{非空} $E_n$，则每个 $E_n$ 都\Term{nonempty}{非空}（因为由映射 $f_n:E_{n+1}\to E_n$ 的存在，$E_n=\varnothing$ 蕴涵 $E_{n+1}=\varnothing$）。称满足 $f_j(x_{j+1})=x_j$（$j=1,\ldots,n-1$）的\Term{finite}{有限}序列 $(x_j)_{1\leq j\leq n}$ 为“相容”的。对 $n$ 归纳证明：存在相容序列 $(a_1,\ldots,a_n)$（$a_i\in E_i$），使对每个 $q\geq n$，它都能延长为长为 $q$ 的相容序列 $(a_1,\ldots,a_n,x_{n+1},\ldots,x_q)$。$n=0$ 时显然，因为每个 $E_q$ 都\Term{nonempty}{非空}。由 $n$ 过渡到 $n+1$。若对每个 $x\in f_n^{-1}(\{a_n\})\subset E_{n+1}$，延长 $(a_1,\ldots,a_n,x)$ 的所有相容序列的长度都以某个整数 $q(x)$ 为界，则因 $E_{n+1}$ \Term{finite}{有限}，延长 $(a_1,\ldots,a_n)$ 的所有相容序列的长度也都有界（以 $\sup_xq(x)$ 为界）；所以存在 $a_{n+1}\in f_n^{-1}(\{a_n\})$，使序列 $(a_1,\ldots,a_n,a_{n+1})$ 能有任意长的延长。于是得到无限序列 $(a_n)_{n\geq1}$，与假设矛盾。
\end{proof}
由引理 3，若 $A$ 是\Term{finite-type}{有限型} \Term{artinian-set}{Artin 集合}，则它的链所成的\Term{ordered-set}{有序集} $G(A)$ 是\Term{finite}{有限}的：事实上，以 $E_n$ 表示由
\SourcePage{127}
$n+1$ 项组成的链 $(x_n\in x_{n-1}\in\cdots\in x_1\in A)$ 所成的\Term{set}{集合}（它是\Term{finite}{有限}的，因为 $A$ 的阶数 $\leq n+1$ 的\Term{set-component}{分量}只有\Term{finite}{有限}多个，且全都是\Term{finite}{有限}集），以 $f_n$ 表示映射
\[
 (x_{n+1}\in x_n\in x_{n+1}\cdots)\longmapsto(x_n\in x_{n-1}\in\cdots).
\]
而\Term{finite}{有限}\Term{ordered-set}{有序集}的\Term{isomorphism-class}{同构类}所成的\Term{set}{集合}是可数的：事实上，在 $\mathbf N$ 的一个\Term{finite}{有限}\Term{subset}{子集}上给定序结构，\Term{equivalence}{等价}于给定它的图，而这个图是 $\mathbf N\times\mathbf N$ 的一个\Term{finite}{有限}\Term{subset}{子集}；另一方面，可数集的所有\Term{finite}{有限}\Term{subset}{子集}所成的\Term{set}{集合}是可数的（第三章）。由 (*) 及引理 1，\Term{finite-type}{有限型} \Term{artinian-set}{Artin 集合}所成的\Term{set}{集合} $\mathfrak a$ 是可数的。由引理 2，它是无限的（或者更简单地，使用由 $z_\circ=\varnothing$、$z_{n+1}=\{z_n\}$ 定义的序列 $(z_n)_{n\geq0}$）。这就完成了 b) 的证明。

\textbf{注记。}我们刚刚证明了：若 $A$ 是\Term{finite-type}{有限型} \Term{artinian-set}{Artin 集合}，则它只有\Term{finite}{有限}多个\Term{set-component}{分量}。因此存在整数 $n$，使 $A$ 为 $n$ 型。

再证明 c)。设 $\mathfrak{c}$ 是无限\Term{cardinal}{基数}，$U$ 是有一个\Term{cardinal}{基数}为 $\mathfrak{c}$ 的元素的\Term{universe}{宇宙}，$A$ 是 $\mathfrak{c}$ 型 \Term{artinian-set}{Artin 集合}。则\Term{ordered-set}{有序集} $G(A)$ 的\Term{cardinal}{基数} $\leq\mathfrak{c}$（上面的公式 (2)）。对 $G(A)$ 应用引理 1 的第二个断言，说明 $G(A)$ 上的映射 $(x_0,\ldots,x_n)\mapsto x_n$ 的值落在 $U$ 中。换言之，$A$ 的所有\Term{set-component}{分量}都是 $U$ 的元素。这就证明了 c)。c$'$) 的证明类似：若 $A$ 是\Term{finite-type}{有限型} \Term{artinian-set}{Artin 集合}，我们已经看到 $G(A)$ \Term{finite}{有限}；因 $U$ \Term{nonempty}{非空}，由第 1 节命题 7 的推论，它包含一个与 $G(A)$ \Term{equipotent}{等势}的元素。证毕。

\SGAsection{7}{几条含混的元数学注记}
a) 公理“每个\Term{set}{集合}都是 \Term{artinian-set}{Artin 集合}”是无害的。事实上，若有\Term{set-theory}{集合论}的一个\Term{model}{模型} $M$，则 $M$ 的所有 \Term{artinian-set}{Artin 元素}所成的\Term{set}{集合} $M'$ 也是一个\Term{model}{模型}（参见命题 11）。

b) \Term{axiom-of-universes}{宇宙公理} (A.6)（以及与它\Term{equivalence}{等价}的\Term{strongly-inaccessible-cardinal}{强不可达基数}公理 (A$'$.6)）独立于\Term{set-theory}{集合论}的其余部分。事实上，设 $\mathfrak{c}$ 是第一个不可数的\Term{strongly-inaccessible-cardinal}{强不可达基数}。严格 $\mathfrak{c}$ 型 \Term{artinian-set}{Artin 集合}所成的\Term{universe}{宇宙} $U_{\mathfrak{c}}$（定理 2 b)）是去掉 (A.6) 的\Term{set-theory}{集合论}的一个\Term{model}{模型}：把 $U_{\mathfrak{c}}$ 的元素称为“\Term{set}{集合}”，“属于关系”是通常的属于关系在 $U_{\mathfrak{c}}$ 上的\Term{restriction}{限制}，等等。因而，\Term{model}{模型}中的“\Term{universe}{宇宙}”就是那些作为 $U_{\mathfrak{c}}$ 的元素的通常\Term{universe}{宇宙}。而我们已经看到，作为 $U_{\mathfrak{c}}$ 的元素的\Term{universe}{宇宙}只有两个小家伙 $U_0=\varnothing$ 和 $U_1$。因此 $U_1$ 是一个不属于任何“\Term{universe}{宇宙}”的“\Term{set}{集合}”。所以，我们有了一个 (A.6) 为假的\Term{set-theory}{集合论}\Term{model}{模型}。

c) Bourbaki 对于\Term{axiom-of-infinity}{无穷公理} (A.5) 独立于前面的公理及公理模式，仅满足于“推测”，未免太谨慎了。它确实是独立的，
\SourcePage{128}
因为由\Term{finite-type}{有限型} \Term{artinian-set}{Artin 集合}组成的可数\Term{universe}{宇宙} $U_1$ 就是一个\Term{model}{模型}，其中 (A.5) 为假，而前面的公理及公理模式都为真。

d) 若能证明\Term{axiom-of-universes}{宇宙公理} (A.6) 是无害的，那将十分有趣。这看来很困难，而且 Paul Cohen 说，它甚至不可证明。

标题中的形容词“含混”，是说我们在构造\Term{model}{模型}时，没有费心把\Term{signature}{符号} $\tau$ 调整妥当，使它不会跑出\Term{model}{模型}。对于\Term{model}{模型} $M$，关键是把通常的 $\tau_x(\mathscr R(x))$ 换成
\[
 \tau_x(\mathscr R(x)\text{ 且 }x\in M).
\]
还须验证，这确实把通常的\Term{quantifier}{量词}变成以前所谓的“典型”\Term{quantifier}{量词}：
\[
 (\forall x\in M)\quad\text{和}\quad(\exists x\in M).
\]

\textbf{练习。}1）设 $n$ 是整数且 $n\geq1$。证明，$n$ 型 \Term{artinian-set}{Artin 集合}所成的\Term{set}{集合}是无限的（\Term{set}{集合} $z_0=\varnothing$、$z_1=\{\varnothing\}$、$z_{q+1}=\{z_q\}$ 都是 $1$ 型的）。

2）称\Term{set}{集合} $A$ 的\Term{set-height}{高度}为满足存在序列 $x_n\in x_{n-1}\in\cdots\in x_0=\square$ 的整数 $n$ 的\Term{supremum}{上确界}（\Term{finite}{有限}或无限）。〔译注：原文等号右端漏印，依本句关于\Term{set}{集合} $A$ 的\Term{height}{高度}的定义及后文递推式补 $A$。〕证明，\Term{set-height}{高度} $\leq n$ 的\Term{set}{集合}都是\Term{finite}{有限}集，且构成一个\Term{finite}{有限}集，它的\Term{cardinal}{基数} $p_n$ 可由 $p_0=1$、$p_{n+1}=2^{p_n}$ 算出（对 $n$ 归纳，注意\Term{set-height}{高度} $\leq n$ 的\Term{set}{集合} $A$ 的元素，都是\Term{set-height}{高度} $\leq n-1$ 的\Term{set}{集合}）。

3）设 $G$ 是 \Term{noetherian-ordered-set}{Noether 有序集}，有最小元素 $g_0$，并且对每个 $h\in G$，所有 $g\leq h$ 所成的\Term{set}{集合}都是\Term{finite-totally-ordered-set}{有限全序集}。

a) 证明 $G$ 的每个元素 $h\ne g_0$ 都有唯一的前驱。

b) 设 $\varphi$ 是引理 1 中定义的 $G$ 上的映射（即 $\varphi(g)$ 是所有 $\varphi(g')$ 所成的\Term{set}{集合}，其中 $g'$ 遍历 $g$ 的各后继）。令 $A=\varphi(g_0)$。证明 $A$ 是 \Term{artinian-set}{Artin 集合}。对 $g\in G$，设 $g_0<g_1<\cdots<g_n=g$ 是所有 $\leq g$ 的元素组成的序列；令 $f(g)=(\varphi(g_0),\varphi(g_1),\ldots,\varphi(g_n))$；证明 $f$ 是从 $G$ 到正文中的\Term{ordered-set}{有序集} $G(A)$ 的递增\Term{surjectivity}{满射}。证明若 $f$ \Term{injectivity}{单射}，则它是 $G$ 到 $G(A)$ 的\Term{isomorphism}{同构}。

c) 对 $g\in G$，设 $M_g$ 是 $g$ 的各上界所成的\Term{set}{集合}。假定对具有同一前驱 $p$ 的任意两个不同元素 $g,g'$，\Term{ordered-set}{有序集} $M_g$ 与 $M_{g'}$ 不\Term{isomorphism}{同构}。证明 b) 中的映射 $f$ 于是为 $G$ 到 $G(A)$ 的\Term{isomorphism}{同构}。[若 $f(g)=f(g')$ 而 $g\ne g'$，并且 $g_0<g_1<\cdots<g_n=g$ 及 $g'_0<g'_1<\cdots<g'_{n'}=g'$ 分别是所有 $\leq g$ 及所有 $\leq g'$ 的元素组成的序列，证明 $n=n'$，且可以假定 $g,g'$ 有同一前驱 $p$；接着考虑满足存在两个不同后继 $g,g'$ 且 $f(g)=f(g')$ 的 $p\in G$ 所成的\Term{set}{集合}，取这个\Term{set}{集合}的一个极大元素 $q$，再取 $q$ 的两个不同后继 $h,h'$，使 $f(h)=f(h')$；注意 $f$ 在 $M_h$ 和 $M_{h'}$ 上的\Term{restriction}{限制}都是\Term{injectivity}{单射}，因而
\SourcePage{129}
（由 b)）分别是 $M_h$ 到 $G(\varphi(h))$ 及 $M_{h'}$ 到 $G(\varphi(h'))$ 的\Term{isomorphism}{同构}；由 $M_h,M_{h'}$ 不\Term{isomorphism}{同构}的假设，推出 $\varphi(h)\ne\varphi(h')$，这与 $f(h)=f(h')$ 矛盾。]

\textbf{附记。}练习 3）十分精确地说明了 \Term{artinian-set}{Artin 集合}是怎样“制造”出来的。我们已经知道，这样的\Term{set}{集合} $A$ 由\Term{ordered-set}{有序集} $G(A)$ 的\Term{isomorphism-class}{同构类}确定（引理 1）。现在可以刻画与某个 $G(A)$ \Term{isomorphism}{同构}的\Term{ordered-set}{有序集} $G$：

$\alpha$）$G$ 是 \Term{noetherian-ordered-set}{Noether 的}，且有最小元素；

$\beta$）对每个 $g\in G$，所有 $h\leq g$ 所成的\Term{set}{集合}都是\Term{finite-totally-ordered-set}{有限全序集}（因而 $g$ 的前驱存在且唯一）；

$\gamma$）若 $g,g'$ 是具有同一前驱的两个不同元素，则 $g$ 的上界所成的\Term{set}{集合} $M_g$ 与 $g$ 的上界所成的\Term{set}{集合} $M_{g'}$ 不\Term{isomorphism}{同构}。

4）设 $(A_n)_{n\geq0}$ 是\Term{finite-ordinal}{有限序数}所成的序列，定义为 $A_0=\varnothing$、$A_{n+1}=A_n\cup\{A_n\}$。证明\Term{ordered-set}{有序集} $G(A_n)$ 有 $2^n$ 个元素，其中\Term{set-height}{高度}为 $q$（按练习 2）的意义）的元素个数为 $\binom nq$。

\SGAsection{}{参考文献}
\begin{enumerate}
\item[{[1]}] B. Mitchell，\textit{\Term{category}{范畴}论}（\textit{Theory of categories}），Academic Press，1965。
\end{enumerate}
\SourcePage{130}
% 原 PDF 第 130 页为空白页。
